\section{Implementation History and Codebase Overview}
\label{sec:history}

\subsection{Version history}

The repository contains 30 commits from 23 June 2026 to 8 October 2026, with two contributors. It
has no tags, and the version history was built from commit comments and dates rather than from
release markers. Table~\ref{tab:history} summarises the seven phases in the order they occurred,
with the labels used in the project's presentation material (the informal version).

\begin{table}[H]
\centering
\caption{Development phases reconstructed from the commit history.}
\label{tab:history}
\footnotesize
\begin{tabular}{@{}L{1.3cm}L{2.0cm}L{8.2cm}L{2.6cm}@{}}
\toprule
\textbf{Phase} & \textbf{Dates} & \textbf{Content} & \textbf{Deck label} \\
\midrule
1 & 23 Jun 2026 & Repository bootstrapped; first generators; first exported solutions. & v1 \\
2 & 29 Jul -- 22 Aug 2026 & Preference-based enrolment simulation; time-pattern update; elective
    lecture/lab pairing fix on the course offering import. & v2, v3, v5 \\
3 & 26--29 Aug 2026 & Constraint work: term advanced to Spr5, then to Spr6; minor-degree block
    scheduling; flexible time-preference rule engine; emission of matching slots, instead of
    blanket prohibitions; master pipeline \texttt{gen\_all.py}; lunch-break persistence repair;
    purge utility; renamed the first-year reserved course. & v4 \\
4 & 7--11 Sep 2026 & Full Taasika dump is added; snapshot 263 is integrated into the loader, with
    improvements to the KPI evaluation, and a room-preference toggle has been added. & --- \\
5 & 28 Sep 2026 & The reinforcement-learning module and the solution for Spring 2026 land directly
    in one commit, along with the setup guide via arm B1. & RL \\
6 & 29 Sep -- 1 Oct 2026 & Departmental and open-elective logic and curriculum correction; open
    elective that can only appear in second year; automatic archive generation; online rooms
    removed from generation; open-elective block enrolments added. & --- \\
7 & 8 Oct 2026 & Student group membership is based on enrolments. & --- \\
\bottomrule
\end{tabular}
\end{table}

The elective pairing fix is recorded in the handover notes as an offering reduction from 63 to 58:
five standalone laboratory offerings were absorbed into their lecture partners. An earlier note
said 63 reduced to 54. That 54 is not in the notes or in the files. The current dataset, after
the later curriculum edits, contains 56 offerings (Table~\ref{tab:dataset}). The phase structure
above is not necessarily in sync with the labels v1--v5, and is only loosely tied to the
individual commits.

\subsection{Repository layout}

\begin{table}[H]
\centering
\caption{Top-level repository structure.}
\label{tab:layout}
\footnotesize
\begin{tabular}{@{}L{3.0cm}L{11.9cm}@{}}
\toprule
\textbf{Path} & \textbf{Contents} \\
\midrule
\texttt{scripts/}     & 25 Python generators, validators and evaluators, and two Jython scripts in
                        \texttt{unitime\_scripts/} that execute inside UniTime \\
\texttt{unitime-out/} & Created XML for the generalised Data Exchange, the frozen dataset
                        \texttt{0gen\_all/} and the archive \\
\texttt{rl/}          & Maven module, Python training package, configuration, instance, models and
                        experiment runs for \texttt{unitime-rl} \\
\texttt{cpsolver/}    & Local copy of the CPSolver library, ready to be built and installed into
                        the module system \\
\texttt{solutions/}   & Exported timetables used on report sheets for the KPI reports to score \\
\texttt{taasika-db/}  & The institutional SQL dumps \\
\texttt{compare/}     & Snapshot summaries usable by humans to compare snapshots \\
\bottomrule
\end{tabular}
\end{table}

Figure~\ref{fig:arch} shows how these parts relate at run time.

\begin{figure}[H]
\centering
\begin{tikzpicture}[
  font=\scriptsize,
  grp/.style={draw,rounded corners=3pt,inner sep=5pt},
  blk/.style={draw,rounded corners=2pt,align=center,minimum height=7mm,inner xsep=4pt},
  ar/.style={-{Stealth[length=1.8mm]},gray!75}
]
% data layer
\node[blk,fill=gray!14] (sql) {\texttt{taasika-db/*.sql}};
\node[blk,fill=gray!14,below=5mm of sql] (snap) {snapshot 242};

% generation layer
\node[blk,fill=coepblue!10,right=12mm of sql,text width=26mm] (scripts) {\texttt{scripts/}\\8 generators\\3 validators};
\node[blk,fill=coepblue!18,below=5mm of scripts,text width=26mm] (xml) {\texttt{unitime-out/0gen\_all}\\16 XML + archive};

% unitime
\node[blk,fill=orange!16,right=12mm of scripts,text width=26mm] (ut) {UniTime server\\{\tiny Java 21, Tomcat 9, MySQL 8.4}};
\node[blk,fill=orange!10,below=5mm of ut,text width=26mm] (cps) {CPSolver\\{\tiny IFS + CBS}};

% rl
\node[blk,fill=green!12,below=12mm of cps,text width=30mm] (rl) {\texttt{rl/} \texttt{unitime-rl}\\{\tiny headless solver, RL layer}};

% outputs
\node[blk,fill=red!8,left=12mm of rl,text width=26mm] (sol) {\texttt{solutions/}\\CSV + solver XML};
\node[blk,fill=red!8,left=12mm of sol,text width=26mm] (kpi) {KPI evaluator\\{\tiny JSON report}};

\draw[ar] (sql) -- (snap);
\draw[ar] (snap.east) -- ++(4mm,0) |- (scripts.west);
\draw[ar] (scripts) -- (xml);
\draw[ar] (xml.east) -- ++(4mm,0) |- (ut.west);
\draw[ar] (ut) -- (cps);
\draw[ar] (cps) -- (sol.north east);
\draw[ar] (cps) -- (rl);
\draw[ar] (rl) -- (sol);
\draw[ar] (sol) -- (kpi);
\draw[ar] (xml.south) |- ++(-4mm,-6mm) |- (kpi.north);
\end{tikzpicture}
\caption{The \texttt{scripts/} generators and validators read the Taasika SQL dump (snapshot 242)
and create the 16 XML files and archive in \texttt{unitime-out/0gen\_all}. That output goes to the
UniTime server (Java 21, Tomcat 9, MySQL 8.4) using CPSolver (IFS + CBS), and also to the headless
\texttt{unitime-rl} module with the RL layer. Both solve one XML and write CSV and solver XML to
\texttt{solutions/}. The KPI evaluator scores those solutions against the same XML and returns a
JSON report.}
\label{fig:arch}
\end{figure}

\subsection{The generation and solving procedure}

The procedure actually followed is the one below, in the style of the corresponding section of the
Semester~VI report but updated for the current tooling.

\begin{enumerate}
  \item \textbf{Generate.} \texttt{python scripts/gen\_all.py --term Spr12 --sql
    taasika-db/taasika3.sql --snapshot 242}, the same command as
    Listing~\ref{lst:genall}. The eight generators are dependent on one another.
  \item \textbf{Validate and archive.} The same command runs the three validators and writes
    \texttt{0gen\_all.zip} if all three pass.
  \item \textbf{Import.} Upload the archive through Data Exchange; it uploads files 1--16 one by
    one.
  \item \textbf{Attach buildings and teachers.} Run \textbf{Buildings\,$\rightarrow$\,Update
    Data}. There is no need to send any data at this point merely because a building has no room,
    or each building reports that it does not have room. Add staff through \textbf{Manage
    Instructor List}.
  \item \textbf{Lunch breaks.} Run \texttt{setup\_lunch\_breaks.py} under
    \textbf{Utilities\,$\rightarrow$\,Scripts}.
  \item \textbf{Pin enrolments.} Set up the three functions of Section~\ref{sec:setup}.
  \item \textbf{Check, solve, export.} Run with the \texttt{Check} configuration. It cycles through
    all classes and prints the loaded classes and the number of students loaded (there must be at
    least one). Load in \texttt{Initial} mode, then solve and export CSV until the timeout occurs.
    If the solver XML is to be exported, set \texttt{unitime.solution.export.names=true}
    beforehand.
  \item \textbf{Evaluate and refine.} Repeat the steps above for the same dataset, and score with
    \texttt{eval\_timetable\_kpis.py} using the check to be fixed.
\end{enumerate}

\subsection{The most recent change: enrolment-derived group membership}

In the last commit of the phase, the way student groups are populated was changed. In the past the
data had been reported as 242 groups, back-filled as empty groups, since bulk updates of the
groups caused a Hibernate failure while importing the data. The generator now operates on
enrolments: membership of a class is determined by the schedule note for that class, and this is
emitted as a group reference for that class within the student's record in file 14. If a note is
not recognised it is dropped and not imported. The codes are checked against the codes declared in
file 6.

The result is 5{,}767 group references from 1{,}130 students, verified exhaustively: each student
enrolled in a class is part of the group of that class, for all classes, with no unrecognised code
and well-formed output. As Section~\ref{sec:constraints} explains, this is possible without
altering the timetable, so the cohort structure is visible. With the pinning flags applied, 40
conflicts were present after the change.
